\subsubsection{Subproceso 2: Selección y Registro de Cursos}

A partir de la revisión documental y la observación directa se obtuvieron los siguientes resultados sobre el funcionamiento actual del Subproceso 2:

\textbf{Actividades observadas y flujo efectivo:}
\begin{itemize}
    \item El portal DIRCE muestra a cada estudiante solo los cursos habilitados según su plan de estudios y los prerrequisitos que ya aprobó.
    \item El control de límite de créditos y de cruces de horario funciona al instante: el sistema impide agregar un curso apenas detecta un incumplimiento, sin esperar a que el alumno envíe su horario final.
    \item Al confirmar el envío del horario, el sistema separa de inmediato las vacantes elegidas y ya no permite que el estudiante haga cambios por su cuenta sin acudir a una rectificación.
\end{itemize}

\textbf{Tiempos de espera identificados:}
\begin{itemize}
    \item La pantalla de selección de cursos demora en cargar durante los primeros minutos en que se abren los turnos, tardando entre 15 y 90 segundos en responder según lo reportado por los alumnos.
    \item Al confirmar el horario en momentos de alta demanda, el sistema puede tardar varios minutos en responder, dejando al estudiante con la duda de si su matrícula se guardó bien.
\end{itemize}

\textbf{Problemas identificados:}
\begin{itemize}
    \item \textbf{Rápido agotamiento de vacantes:} Las secciones de los cursos más pedidos se llenan en los primeros minutos de los turnos de mayor promedio, dejando pocas opciones de horario para los alumnos de turnos siguientes.
    \item \textbf{Lentitud y caídas del portal:} Cuando entran muchos estudiantes al mismo tiempo al abrirse un turno, la Intranet DIRCE se vuelve lenta o se cae, dificultando terminar la matrícula a tiempo.
    \item \textbf{Cruces de horario entre ciclos seguidos:} Hay cursos de ciclos contiguos que tienen el mismo horario y no cuentan con secciones adicionales, impidiendo que el estudiante los lleve juntos.
    \item \textbf{Cursos que no aparecen en el sistema:} Se reportaron asignaturas del plan de estudios que no figuraban disponibles en el portal porque faltó cargarlas al programar los horarios en los Departamentos Académicos.
\end{itemize}
